\chapter{Supplement: Sampling Distributions and Descriptive Statistics}

\begin{nuance}{place in the course}
This is a supplementary chapter. It is not part of the textbook's Chapter~7. It collects what the textbook states without proof or leaves to earlier chapters:
\begin{itemize}
  \item descriptive statistics;
  \item the sampling distribution of the sample mean;
  \item the exact distributions of $S^2$, $t$ and $F$ for normal samples, with the proofs the textbook only cites.
\end{itemize}
Chapters~1--3 refer to it.
\end{nuance}

Probability starts from a model and asks what data it is likely to produce. Statistics runs the other way: it starts from data and asks which model could have produced it.

The bridge is the \emph{sampling distribution}. A statistic computed from random data is itself random. Its distribution tells us how far the statistic can be trusted.

This chapter does four things:
\begin{itemize}
  \item sets up the language of populations and samples;
  \item summarises data;
  \item derives the exact sampling distributions of the mean and variance of a normal sample;
  \item introduces the chi-square, $t$ and $F$ distributions, which carry almost all of classical inference.
\end{itemize}

\section{Populations, samples and statistics}

\begin{itemize}
  \item A \emph{population} is the full collection of units of interest (all bolts from a factory, all future customers).
  \item A \emph{parameter} is a fixed number describing it, such as the mean $\mu$ or a proportion $p$.
  \item We rarely observe the whole population, so we observe a \emph{sample}.
\end{itemize}

\begin{definition}[Random sample]
A \emph{random sample} of size $n$ from a distribution $F$ is a set $X_1,\dots,X_n$ of independent random variables, each with distribution $F$. We say the $X_i$ are \emph{i.i.d.} (independent and identically distributed). The observed values $x_1,\dots,x_n$ are the \emph{data}.
\end{definition}

\begin{definition}[Statistic, sampling distribution]
A \emph{statistic} is any function $T=T(X_1,\dots,X_n)$ of the sample that does not involve unknown parameters. Its \emph{sampling distribution} is the distribution of the values $T$ would take over repeated samples of the same size from the same population.
\end{definition}

\begin{intuition}
One data set gives one value of $T$. The sampling distribution is the whole cloud of values you would see if you could repeat the experiment many times.
\begin{itemize}
  \item Inference means using the one value you have to work out what the cloud looks like.
  \item A narrow cloud means a reliable statistic.
\end{itemize}
\end{intuition}

\begin{pitfall}{i.i.d. is a modelling assumption, not a fact}
Almost every formula in this book assumes independent, identically distributed observations. These break it:
\begin{itemize}
  \item clustered data (students in the same school);
  \item repeated measurements on the same unit;
  \item time series;
  \item convenience samples.
\end{itemize}
Using i.i.d. formulas on such data usually understates the uncertainty, sometimes by a lot. How the data were collected matters more than any calculation that follows.
\end{pitfall}

\section{Describing a data set}

Before fitting any model, look at the data. The standard summaries follow.

\begin{definition}[Sample mean, variance, standard deviation]
For data $x_1,\dots,x_n$:
\[ \bar{x}=\frac1n\sum_{i=1}^nx_i,\qquad s^2=\frac1{n-1}\sum_{i=1}^n(x_i-\bar{x})^2,\qquad s=\sqrt{s^2}. \]
\end{definition}

The divisor $n-1$ (not $n$) makes $s^2$ an unbiased estimator of the population variance. Chapter~1 (Section 1.2) explains why.

A handy identity:
\[ \sum_{i=1}^n(x_i-\bar{x})^2=\sum_{i=1}^nx_i^2-n\bar{x}^2. \]
For numerical safety, prefer a two-pass or Welford-type algorithm. The identity subtracts two large, nearly equal numbers.

\begin{definition}[Order statistics, median, quantiles]
Sort the data: $x_{(1)}\le x_{(2)}\le\dots\le x_{(n)}$.
\begin{itemize}
  \item The \emph{median} is $x_{((n+1)/2)}$ for odd $n$, and $\tfrac12(x_{(n/2)}+x_{(n/2+1)})$ for even $n$.
  \item The \emph{$p$-th sample quantile} has a fraction $p$ of the data at or below it. (There are several interpolation conventions.)
  \item The quartiles $Q_1,Q_2,Q_3$ sit at $p=0.25,0.5,0.75$.
  \item The \emph{interquartile range} is $\mathrm{IQR}=Q_3-Q_1$.
\end{itemize}
\end{definition}

\begin{example}
Ten shaft diameters (mm): $9.9,\ 10.1,\ 10.0,\ 10.2,\ 9.8,\ 10.0,\ 10.1,\ 9.9,\ 10.0,\ 12.0$.
\begin{itemize}
  \item Mean: $\bar{x}=10.2$. Standard deviation: $s\approx0.64$.
  \item Sorted: $9.8,9.9,9.9,10.0,10.0,10.0,10.1,10.1,10.2,12.0$. Median $=10.0$. IQR $\approx0.2$.
  \item The single outlier $12.0$ pulls the mean up to $10.2$ and inflates $s$. It barely moves the median or the IQR.
\end{itemize}
\end{example}

\begin{theorem}[Robustness comparison]\label{thm:robust}
The mean and standard deviation have \emph{breakdown point} $0$: one wild observation can push them arbitrarily far. The median has breakdown point $1/2$: nearly half the data can be corrupted without moving the median outside the range of the clean data.
\end{theorem}

\begin{proof}
Replace $x_n$ by $t\to\infty$. Then $\bar{x}\to\infty$ and $s\to\infty$. For the median, as long as fewer than $\lceil n/2\rceil$ points are altered, the middle order statistic stays between two clean order statistics.
\end{proof}

\begin{interview}{mean or median?}
\begin{itemize}
  \item Use the mean for symmetric, light-tailed data when efficiency matters.
  \item Use the median and IQR for skewed or contaminated data: incomes, latencies, response times.
  \item Report both when they differ strongly. The gap itself says something about skew.
  \item Related fact: the mean minimises $\sum(x_i-c)^2$ over $c$. The median minimises $\sum\abs{x_i-c}$.
\end{itemize}
\end{interview}

\subsection*{Graphical summaries}

\begin{itemize}
  \item A \emph{histogram} splits the range into bins and counts the data in each. Its shape shows skew, several modes and outliers.
  \item A \emph{box plot} shows the median, the quartiles, and points beyond $1.5\,\mathrm{IQR}$ from the box as possible outliers.
  \item A \emph{normal quantile (Q--Q) plot} graphs the ordered data $x_{(i)}$ against normal quantiles $\Phi^{-1}\big(\tfrac{i-0.5}{n}\big)$. Roughly normal data lie near a straight line. Systematic curvature means skew or heavy tails.
\end{itemize}

\begin{figure}[h]
\centering
\begin{tikzpicture}[scale=1.0]
  \draw[axis] (-3.4,0)--(3.6,0); \draw[axis] (0,-0.2)--(0,2.6);
  \draw[thick,accent2,domain=-3.2:3.2,samples=60] plot (\x,{1.2*exp(-\x*\x/2)});
  \draw[thick,accent,domain=-3.2:3.2,samples=80] plot (\x,{2.4*exp(-2*\x*\x)});
  \node[lab,accent2] at (2.3,0.9){one observation}; \node[lab,accent] at (-1.7,2.2){$\bar{X}$, four observations};
  \node[lab] at (0,-0.6){the distribution of $\bar{X}$ concentrates as $n$ grows};
\end{tikzpicture}
\caption{Two bell curves centred at the true mean: the spread of one observation, and the narrower spread of the average of several observations.}
\end{figure}

\section{The sampling distribution of the sample mean}

\begin{theorem}[Mean and variance of $\bar{X}$]\label{thm:xbar}
If $X_1,\dots,X_n$ are i.i.d. with mean $\mu$ and variance $\sigma^2$, then
\[ \Exp[\bar{X}]=\mu,\qquad\Var(\bar{X})=\frac{\sigma^2}n,\qquad\se(\bar{X})=\frac{\sigma}{\sqrt n}. \]
\end{theorem}

\begin{proof}
By linearity, $\Exp[\bar{X}]=\frac1n\sum\mu=\mu$. By independence the variances add: $\Var(\bar{X})=\frac1{n^2}\sum\sigma^2=\sigma^2/n$.
\end{proof}

The \emph{standard error} of a statistic is its standard deviation, $\se(T)=\sqrt{\Var[T]}$. The $1/\sqrt n$ law is the most important rate in statistics: to halve the uncertainty you need four times the data.

\begin{theorem}[Normal samples]\label{thm:normalmean}
If $X_1,\dots,X_n$ are i.i.d. $\Normal(\mu,\sigma^2)$, then $\bar{X}\sim\Normal(\mu,\sigma^2/n)$ exactly.
\end{theorem}

\begin{proof}
A linear combination of independent normals is normal. (Moment generating functions multiply: $M_{\bar{X}}(t)=\prod M_{X_i}(t/n)=e^{\mu t+\sigma^2t^2/(2n)}$.) The mean and variance are as computed above.
\end{proof}

Without normality, the central limit theorem gives an approximation.

\begin{theorem}[CLT for the sample mean]\label{thm:clt}
If the $X_i$ are i.i.d. with mean $\mu$ and finite variance $\sigma^2>0$, then $\sqrt n(\bar{X}-\mu)/\sigma\to\Normal(0,1)$ in distribution as $n\to\infty$.
\end{theorem}

The CLT is proved in the Probability book (Limit Theorems). In practice:
\begin{itemize}
  \item For moderately skewed data, $n\ge30$ is a common rule of thumb.
  \item For heavily skewed or heavy-tailed data, much larger samples may be needed.
  \item For distributions with no finite variance (the Cauchy), the theorem does not apply. The mean of $n$ Cauchy variables has the same distribution as one.
\end{itemize}

\section{The chi-square distribution and the sample variance}

\begin{definition}[Chi-square distribution]
If $Z_1,\dots,Z_\nu$ are i.i.d. $\Normal(0,1)$, then $V=Z_1^2+\dots+Z_\nu^2$ has the \emph{chi-square distribution with $\nu$ degrees of freedom}, $V\sim\chi^2_\nu$.

It is the gamma distribution with shape $\nu/2$ and scale $2$. Its density is $f(v)=\dfrac{v^{\nu/2-1}e^{-v/2}}{2^{\nu/2}\Gamma(\nu/2)}$ for $v>0$.
\end{definition}

\begin{proposition}[Properties of $\chi^2_\nu$]\label{prop:chisq}
$\Exp[V]=\nu$ and $\Var[V]=2\nu$. If $V_1\sim\chi^2_{\nu_1}$ and $V_2\sim\chi^2_{\nu_2}$ are independent, then $V_1+V_2\sim\chi^2_{\nu_1+\nu_2}$.
\end{proposition}

\begin{proof}
$\Exp[Z^2]=1$ and $\Var[Z^2]=\Exp[Z^4]-1=3-1=2$. Then $\nu$ independent terms add their means and variances. Additivity is clear from the definition as a sum of squares (or from the gamma moment generating function $(1-2t)^{-\nu/2}$).
\end{proof}

The next result is the cornerstone of exact normal-theory inference.

\begin{theorem}[Distribution of $\bar{X}$ and $S^2$ for a normal sample]\label{thm:cochran}
Let $X_1,\dots,X_n$ be i.i.d. $\Normal(\mu,\sigma^2)$, with $S^2=\frac1{n-1}\sum(X_i-\bar{X})^2$. Then
\begin{enumerate}[label=(\roman*)]
  \item $\bar{X}$ and $S^2$ are independent;
  \item $\dfrac{(n-1)S^2}{\sigma^2}\sim\chi^2_{n-1}$.
\end{enumerate}
\end{theorem}

\begin{proof}
Standardise: $Z_i=(X_i-\mu)/\sigma$ are i.i.d. $\Normal(0,1)$, so $\mathbf{Z}=(Z_1,\dots,Z_n)\T\sim\Normal(\mathbf{0},I_n)$.

Choose an orthogonal matrix $Q$ whose first row is $\mathbf{1}\T/\sqrt n$. (Extend that unit vector to an orthonormal basis; see Chapter~11 of the Linear Algebra book.) Put $\mathbf{Y}=Q\mathbf{Z}$. Since $Q$ is orthogonal, $\mathbf{Y}\sim\Normal(\mathbf{0},QQ\T)=\Normal(\mathbf{0},I_n)$. The $Y_i$ are again i.i.d. standard normal.

The first coordinate is $Y_1=\sqrt n\,\bar{Z}$. Orthogonal maps preserve length, so $\sum Y_i^2=\sum Z_i^2$. Therefore
\[ \sum_{i=1}^n(Z_i-\bar{Z})^2=\sum Z_i^2-n\bar{Z}^2=\sum_{i=1}^nY_i^2-Y_1^2=\sum_{i=2}^nY_i^2. \]
This is a sum of $n-1$ squared independent standard normals, so it is $\chi^2_{n-1}$. It depends only on $Y_2,\dots,Y_n$, while $\bar{Z}=Y_1/\sqrt n$ depends only on $Y_1$. So the two are independent.

Finally, $(n-1)S^2/\sigma^2=\sum(Z_i-\bar{Z})^2$ and $\bar{X}=\mu+\sigma\bar{Z}$.
\end{proof}

\begin{corollary}[Mean and variance of $S^2$]\label{cor:s2}
For normal samples, $\Exp[S^2]=\sigma^2$ and $\Var[S^2]=2\sigma^4/(n-1)$. (The unbiasedness $\Exp[S^2]=\sigma^2$ holds for every distribution with finite variance; see Chapter~1.)
\end{corollary}

\begin{proof}
$\Exp[\chi^2_{n-1}]=n-1$ and $\Var[\chi^2_{n-1}]=2(n-1)$. Then $S^2=\sigma^2V/(n-1)$ with $V\sim\chi^2_{n-1}$.
\end{proof}

\begin{nuance}{why $n-1$ degrees of freedom}
The $n$ deviations $X_i-\bar{X}$ must sum to zero. So only $n-1$ of them are free. They lie in an $(n-1)$-dimensional subspace: the orthogonal complement of $\mathbf{1}$.

That is what ``degrees of freedom'' means: the dimension of the subspace where the sum of squares lives. The same count reappears in regression and ANOVA (Chapters~5--8). Estimating $p$ parameters removes $p$ degrees of freedom.
\end{nuance}

\section{The $t$ and $F$ distributions}

The standardised mean $(\bar{X}-\mu)/(\sigma/\sqrt n)$ is standard normal. In practice $\sigma$ is unknown, and we replace it by $S$. The distribution changes.

\begin{definition}[Student's $t$ distribution]
If $Z\sim\Normal(0,1)$ and $V\sim\chi^2_\nu$ are independent, then $T=Z/\sqrt{V/\nu}$ has the \emph{$t$ distribution with $\nu$ degrees of freedom}, $T\sim t_\nu$.
\end{definition}

\begin{theorem}[The $t$ statistic]\label{thm:tstat}
For a normal sample, $T=\dfrac{\bar{X}-\mu}{S/\sqrt n}\sim t_{n-1}$.
\end{theorem}

\begin{proof}
Write
\[ T=\dfrac{(\bar{X}-\mu)/(\sigma/\sqrt n)}{\sqrt{S^2/\sigma^2}}=\dfrac{Z}{\sqrt{V/(n-1)}}, \]
with $Z=\sqrt n(\bar{X}-\mu)/\sigma\sim\Normal(0,1)$ and $V=(n-1)S^2/\sigma^2\sim\chi^2_{n-1}$. They are independent by Theorem~\ref{thm:cochran}(i).
\end{proof}

Properties of $t_\nu$:
\begin{itemize}
  \item Symmetric about $0$, with density $f(t)=\dfrac{\Gamma(\frac{\nu+1}2)}{\sqrt{\nu\pi}\,\Gamma(\frac\nu2)}\big(1+\tfrac{t^2}\nu\big)^{-(\nu+1)/2}$.
  \item Heavier tails than the normal. This is the price of estimating $\sigma$.
  \item $\Var[T]=\nu/(\nu-2)$ for $\nu>2$.
  \item $t_1$ is the Cauchy distribution.
  \item $t_\nu\to\Normal(0,1)$ as $\nu\to\infty$, because $V/\nu\to1$ by the law of large numbers.
\end{itemize}

\begin{figure}[h]
\centering
\begin{tikzpicture}[scale=1.0]
  \draw[axis] (-4.4,0)--(4.6,0); \draw[axis] (0,-0.2)--(0,2.8);
  \draw[thick,accent2,domain=-4:4,samples=70] plot (\x,{2.6*exp(-\x*\x/2)});
  \draw[thick,accent,domain=-4:4,samples=70] plot (\x,{2.6*0.9*(1+\x*\x/3)^(-2)});
  \node[lab,accent2] at (1.9,1.5){$\Normal(0,1)$}; \node[lab,accent] at (2.4,0.5){$t_3$: heavier tails};
\end{tikzpicture}
\caption{The $t_3$ density has a lower peak and heavier tails than the standard normal. Large observations of $T$ are much more likely when $\sigma$ must be estimated from a small sample.}
\end{figure}

\begin{definition}[$F$ distribution]
If $U\sim\chi^2_{a}$ and $V\sim\chi^2_{b}$ are independent, then $F=\dfrac{U/a}{V/b}$ has the \emph{$F$ distribution with $(a,b)$ degrees of freedom}, $F\sim F_{a,b}$.
\end{definition}

Two quick consequences:
\begin{itemize}
  \item $1/F\sim F_{b,a}$.
  \item If $T\sim t_\nu$, then $T^2\sim F_{1,\nu}$, because $T^2=Z^2/(V/\nu)$ with $Z^2\sim\chi^2_1$.
\end{itemize}

\begin{theorem}[Ratio of sample variances]\label{thm:fratio}
Let $X_1,\dots,X_m$ be i.i.d. $\Normal(\mu_1,\sigma_1^2)$, and independently $Y_1,\dots,Y_n$ i.i.d. $\Normal(\mu_2,\sigma_2^2)$, with sample variances $S_1^2,S_2^2$. Then
\[ \frac{S_1^2/\sigma_1^2}{S_2^2/\sigma_2^2}\sim F_{m-1,\,n-1}. \]
\end{theorem}

\begin{proof}
By Theorem~\ref{thm:cochran}, $(m-1)S_1^2/\sigma_1^2\sim\chi^2_{m-1}$ and $(n-1)S_2^2/\sigma_2^2\sim\chi^2_{n-1}$. They are independent because the samples are. Divide each by its degrees of freedom.
\end{proof}

\begin{interview}{why is normality so central here?}
Theorems~\ref{thm:cochran}--\ref{thm:fratio} are \emph{exact} only for normal data. In fact, independence of $\bar{X}$ and $S^2$ \emph{characterises} the normal distribution.
\begin{itemize}
  \item For non-normal data, the results for the mean become large-sample approximations through the CLT.
  \item The $\chi^2$ and $F$ results for variances are \emph{not} robust to non-normality, because they depend on the fourth moment.
  \item When normality is doubtful, use the bootstrap or permutation methods (Chapter~13).
\end{itemize}
\end{interview}

\section{Sampling distributions in practice}

\begin{example}[Using the $t$ distribution]
A sample of $n=9$ from a normal population has $\bar{x}=52$ and $s=6$.
\begin{itemize}
  \item Standard error: $s/\sqrt n=2$. The statistic $T=(\bar{X}-\mu)/2$ follows $t_8$.
  \item $t_{8,0.025}\approx2.306$. So with probability $0.95$, $\abs{\bar{X}-\mu}\le2.306\times2=4.61$.
  \item The normal approximation would give $1.96\times2=3.92$. Ignoring the $t$ correction would understate the uncertainty by about $18\%$.
\end{itemize}
\end{example}

\begin{example}[Probability about a sample variance]
A normal population has $\sigma^2=4$. For a sample of $n=10$:
\[ \Prob(S^2>6)=\Prob\big(\chi^2_9>9\cdot6/4=13.5\big)\approx0.14. \]
\end{example}

\begin{summarybox}
\begin{itemize}
  \item A statistic is a function of the sample. Its sampling distribution (what it does over repeated samples) is the basis of all inference.
  \item Descriptive statistics: mean, variance, median, quantiles, IQR, histograms, box plots, Q--Q plots. The median and IQR are robust. The mean and variance are not.
  \item $\Exp[\bar{X}]=\mu$ and $\Var[\bar{X}]=\sigma^2/n$. For normal data $\bar{X}\sim\Normal(\mu,\sigma^2/n)$ exactly, and approximately for large $n$ by the CLT.
  \item For a normal sample, $\bar{X}$ and $S^2$ are independent and $(n-1)S^2/\sigma^2\sim\chi^2_{n-1}$. The proof is an orthogonal change of variables.
  \item $t=Z/\sqrt{V/\nu}$ appears when $\sigma$ is estimated: $(\bar{X}-\mu)/(S/\sqrt n)\sim t_{n-1}$. The ratio of independent scaled $\chi^2$ variables is $F$, which governs variance comparisons.
\end{itemize}
\end{summarybox}

\section{Exercises}

\begin{exercise}
For the data $3,5,5,7,8,9,12$ compute the mean, median, variance, standard deviation, quartiles and IQR. What happens to each if $12$ is replaced by $120$?
\end{exercise}

\begin{exercise}
Prove that $\sum(x_i-\bar{x})=0$ and that the mean minimises $\sum(x_i-c)^2$. Prove that the median minimises $\sum\abs{x_i-c}$ for $n$ odd.
\end{exercise}

\begin{exercise}
Show that if $Y=aX+b$ then $\bar{y}=a\bar{x}+b$ and $s_y=\abs{a}s_x$. What do the quartiles and the IQR do under this transformation?
\end{exercise}

\begin{exercise}
A machine fills bottles with mean $500$ ml and standard deviation $4$ ml. What is the approximate probability that the mean of $64$ bottles exceeds $501$ ml? What assumptions are used?
\end{exercise}

\begin{exercise}
If $Z_1,Z_2,Z_3$ are i.i.d. standard normal, find the distribution of $(Z_1-Z_2)^2/2$ and of $(Z_1+Z_2+Z_3)^2/3$. Are they independent? Justify.
\end{exercise}

\begin{exercise}
For a sample of size $n=16$ from $\Normal(\mu,\sigma^2)$, find $c$ so that $\Prob(S^2\le c\sigma^2)=0.95$ (use $\chi^2_{15}$ quantiles from software).
\end{exercise}

\begin{exercise}
Prove that $T\sim t_\nu$ has $\Exp[T]=0$ for $\nu>1$ and $\Var[T]=\nu/(\nu-2)$ for $\nu>2$ by writing $T=Z/\sqrt{V/\nu}$ and using $\Exp(1/V)=1/(\nu-2)$.
\end{exercise}

\begin{exercise}
Simulate $10{,}000$ samples of size $5$ from a normal population and from an exponential population; compute the $t$ statistic for each and compare histograms with the $t_4$ density. What do you conclude about the effect of non-normality at small $n$?
\end{exercise}
